
In this section, we develop the DFIV algorithm. Similarly to \citet{Singh2019}, we assume that we do not necessarily have  access to observations from the joint distribution of $(X,Y,Z)$. Instead, we are given $m$ observations of $(X,Z)$ for stage 1 and $n$ observations of $(Y,Z)$ for stage 2. We denote the stage 1 observations as $(x_i, z_i)$ and the stage 2 observations as $(\tilde y_i, \tilde z_i)$. If observations of $(X,Y,Z)$ are given for both stages, we can evaluate the out-of-sample losses, and these losses can be used for hyper-parameter tuning of $\lambda_1,\lambda_2$ (Appendix~\ref{sec:hyper-param}).

DFIV uses the following models
\begin{align}
f_\mathrm{struct}(x) = \vec{u}^\top \vec{\psi}_{\theta_X}(x) \hspace{.5cm} \mathrm{and} \hspace{.5cm} \expect[X|z]{\vec{\psi}_{\theta_X}(X)} = \vec{V}\vec{\phi}_{\theta_Z}(z),
\end{align}
where $\vec{u} \in \mathbb{R}^{d_1}$ and $\vec{V} \in \mathbb{R}^{d_1\times d_2}$ are the parameters, and $\vec{\psi}_{\theta_X}(x) \in \mathbb{R}^{d_1}$ and $\vec{\phi}_{\theta_Z}(z)\in \mathbb{R}^{d_2}$ are the neural nets parameterised by $\theta_X \in \Theta_X$ and $\theta_Z \in \Theta_Z$, respectively. As in the original 2SLS algorithm, we learn $\expect[X|z]{\vec{\psi}_{\theta_X}(X)}$ in stage 1 and $f_\mathrm{struct}(x)$ in stage 2. In addition to the weights $\vec{u}$ and $\vec{V}$, however, we also learn the parameters of the feature maps, $\theta_X$ and $\theta_Z$. Hence, we need to alternate between stages 1 and 2, since the conditional expectation $\expect[X|z]{\vec{\psi}_{\theta_X}(X)}$ changes during training.

\paragraph{Stage 1 Regression}
The goal of stage 1 is to estimate the conditional expectation $\expect[X|z]{\vec{\psi
}_{\theta_X}(X)} \simeq \vec{V} \vec{\phi}_{\theta_Z}(z)$ by learning the matrix $\vec{V}$ and parameter $\theta_Z$, with $\theta_X=\hat{\theta}_X$ given and fixed. Given the stage 1 data $(x_i, z_i)$, this can be done by minimizing the empirical estimate of $\mathcal{L}_1$,
\begin{align}
    \hat{\vec{V}}^{(m)}, \hat{\theta}_Z = \argmin_{\vec{V} \in \mathbb{R}^{d_1\times d_2}, \theta_Z \in \Theta_Z} \mathcal{L}_1^{(m)}(\vec{V}, \theta_Z),~ \mathcal{L}_1^{(m)} = \frac1m \sum_{i=1}^m \|\vec{\psi}_{\hat\theta_X}(x_i) - \vec{V}\vec{\phi}_{\theta_Z}(z_i)\|^2 + \lambda_1 \|\vec{V}\|^2. \label{eq:stage1-empirical-loss}
\end{align}
Note that the feature map $\vec{\psi}_{\hat\theta_X}(X)$ is fixed during stage 1, since this is the ``target variable.'' 
If we fix $\theta_Z$, the minimization problem \eqref{eq:stage1-empirical-loss} reduces to a  linear ridge regression problem with multiple targets, whose solution as a function of $\theta_X$ and $\theta_Z$ is given analytically by
\begin{align}
    \hat{\vec{V}}^{(m)}(\theta_X, \theta_Z) = \Psi_1^\top \Phi_1 (\Phi_1^\top \Phi_1 + m\lambda_1 I)^{-1}, \label{eq:stage1-sol}
\end{align}
where $\Phi_1, \Psi_1$ are feature matrices defined as
$
    \Psi_1 = [\vec{\psi}_{\theta_X}(x_1), \dots, \vec{\psi}_{\theta_X}(x_m)]^\top \in \mathbb{R}^{m \times d_1}$ and $
    \Phi_1 = [\vec{\phi}_{\theta_Z}(z_1), \dots, \vec{\phi}_{\theta_Z}(z_m)]^\top \in \mathbb{R}^{m \times d_2}.
$
We can then learn the parameters $\theta_Z$ of the adaptive features $\vec{\psi}_{\theta_Z}$ by minimizing the loss $\mathcal{L}^{(m)}_1$ at $\vec{V} = \hat{\vec{V}}^{(m)}(\hat\theta_X, \theta_Z)$ using gradient descent.
For simplicity, we introduce a small abuse of notation by denoting as $\hat{\theta}_Z$ the result of a user-chosen number of gradient descent steps on the loss \eqref{eq:stage1-empirical-loss} with $\hat{\vec{V}}^{(m)}(\hat\theta_X, \theta_Z)$ from \eqref{eq:stage1-sol}, even though $\hat{\theta}_Z$ need not attain the minimum of the non-convex loss \eqref{eq:stage1-empirical-loss}. We then write $ \hat{\vec{V}}^{(m)}:=\hat{\vec{V}}^{(m)}(\hat\theta_X, \hat{\theta}_Z)$.
While this trick of using an analytical estimate of the linear output weights of a deep neural network might not lead to significant gains in standard supervised learning, it turns out to be very important in the development of our 2SLS algorithm. As shown in the following section, the analytical estimate $\hat{\vec{V}}^{(m)}(\theta_X,\hat{\theta}_Z)$ (now considered as a function of $\theta_X$) will be used to backpropagate to $\theta_X$ in stage 2. 

\paragraph{Stage 2 Regression}
In stage 2, we learn the structural function by computing the weight vector $\vec{u}$ and parameter $\theta_X$ while \emph{fixing}  $\theta_Z=\hat{\theta}_Z$, and thus the corresponding feature map $\vec{\phi}_{\hat{\theta}_Z}(z)$. Given the data $(\tilde y_i, \tilde{z_i})$, we can minimize the empirical version of $\mathcal{L}_2$, defined as
\begin{align}
    \hat{\vec{u}}^{(n)}, \hat{\theta}_X = \argmin_{\vec{u} \in \mathbb{R}^{d_1}, \theta_X \in \Theta_X}\mathcal{L}^{(n)}_2(\vec{u},\theta_X), \quad \mathcal{L}^{(n)}_2 = \frac1n \sum_{i=1}^n (\tilde y_i - \vec{u}^\top \hat{\vec{V}}^{(m)} \vec{\phi}_{\hat\theta_Z}(\tilde z_i))^2 + \lambda_2 \|\vec{u}\|^2. \label{eq:stage2-empirical-loss}
\end{align}

Again, for a given $\theta_X$, we can solve the minimization problem \eqref{eq:stage2-empirical-loss} for $\vec{u}$ as a function of $\hat{\vec{V}}^{(m)}:=\hat{\vec{V}}^{(m)}(\theta_X,\hat\theta_Z)$ by a linear ridge regression
\begin{align}
    \hat{\vec{u}}^{(n)}(\theta_X,\hat\theta_Z) = \left(\hat{\vec{V}}^{(m)}\Phi_2 ^\top \Phi_2 (\hat{\vec{V}}^{(m)})^\top + n\lambda_2 I\right)^{-1}\hat{\vec{V}}^{(m)}\Phi_2^\top \vec{y}_2, \label{eq:stage2-sol}
\end{align}
where
%\begin{align*}
$    \Phi_2 = [\vec{\phi}_{\hat\theta_Z}(\tilde z_1), \dots, \vec{\phi}_{\hat\theta_Z}(\tilde z_n)]^\top \in \mathbb{R}^{n \times d_2}$ and $\vec{y}_2 = [\tilde y_1, \dots, \tilde y_n]^\top \in \mathbb{R}^{n}.
$ %\end{align*}
%Here, $\hat{\vec{u}}^{(n)}$ is a function of $\theta_X$ through the dependence of $\hat{\vec{V}}^{(m)}$. Then, again, we learn $\theta_X$ by minimizing the loss $\mathcal{L}^{(n)}_2$ at $\vec{u} = \hat{\vec{u}}^{(n)}$. 

The loss $\mathcal{L}^{(n)}_2$ explicitly depends on the parameters $\theta_X$ and we can backpropagate it to $\theta_X$ via $\hat{\vec{V}}^{(m)}(\theta_X,\hat{\theta}_Z)$, 
%% AG:I removed the hat above on \theta_Z for consistency with the rest of the section. AD: got it.
%\adcomment{(recall $\theta_Z=\hat{\theta}_Z$ here)},
%AG: do we need this reminder?  I think it's clear AD: ok Arthur
even though the samples of the treatment variable $X$ do not appear in stage 2 regression.
%\adcomment{We again slightly abuse notation and denote by $\hat{\theta}_X$ the parameter estimate obtained after a few gradient steps even if it does not minimize the non-convex loss \eqref{eq:stage2-sol} 
%and write $\hat{\vec{u}}^{(n)}=\hat{\vec{u}}^{(n)}(\hat{\theta}_X,\theta_Z)$.}
We again introduce a small abuse of notation for simplicity, 
and denote by $\hat{\theta}_X$ the  estimate obtained after a few gradient steps on \eqref{eq:stage2-empirical-loss}
with $\hat{\vec{u}}^{(n)}(\theta_X,\hat\theta_Z)$ from \eqref{eq:stage2-sol},
even though $\hat{\theta}_X$ need not minimize the non-convex loss \eqref{eq:stage2-empirical-loss}. 
We then have 
$\hat{\vec{u}}^{(n)}=\hat{\vec{u}}^{(n)}(\hat{\theta}_X,\hat\theta_Z)$.
 %Note that we do not attempt to backpropagate through the estimate $\hat{\theta}_Z$ as this is too computationally expensive. Instead, we rely on alternating optimization of stages 1 and 2. %, and on backpropagation through $\hat{\vec{V}}^{(m)}(\theta_X,\hat{\theta}_Z)$.
%After updating $\hat{\theta}_X$, we need to perform stage 1 regression to update $\hat{\theta}_Z$ again since 
%$\vec{\psi}_{\theta_X}(x)$
% has changed.
After updating $\hat{\theta}_X$, we need to update $\hat{\theta}_Z$ accordingly. We do not attempt to backpropagate through the estimate $\hat{\theta}_Z$ to do this, however, as this would be too computationally expensive; instead, we alternate stages 1 and 2. We also considered updating $\hat{\theta}_X$ and $\hat{\theta}_Z$ jointly to optimize the loss $\mathcal{L}_2^{(n)}$, but this fails, as discussed in Appendix~\ref{sec:joint}.

\paragraph{Computational Complexity and Convergence}

The computational complexity of the algorithm is $O(md_1d_2 + d_2^3)$ for stage 1, while stage 2 requires additional $O(nd_1d_2 + d_1^3)$ computations. This is small compared to KIV \citep{Singh2019}, which takes $O(m^3)$ and $O(n^3)$, respectively. We can further speed up the learning by using mini-batch training as shown in Algorithm~\ref{algo}. Here, $\hat{\vec{V}}^{(m_b)}$ and $\hat{\vec{u}}^{(n_b)}$ are the functions given by \eqref{eq:stage1-sol} and \eqref{eq:stage2-sol} calculated using mini-batches of data. Similarly, $\mathcal{L}_1^{(m_b)}$ and $\mathcal{L}_1^{(n_b)}$ are the stage 1 and 2 losses for the mini-batches. We recommend setting the batch size large enough so that $\hat{\vec{V}}^{(m_b)}, \hat{\vec{u}}^{(n_b)}$ do not diverge from $\hat{\vec{V}}^{(m)}, \hat{\vec{u}}^{(n)}$ computed on the entire dataset. Furthermore, we observe that setting $T_1 > T_2$, i.e. updating $\theta_Z$ more frequently than $\theta_X$, stabilizes the learning process.

 \begin{algorithm}
 \caption{Deep Feature Instrumental Variable Regression}
 \begin{algorithmic}[1]  \label{algo}
 \renewcommand{\algorithmicrequire}{\textbf{Input:}}
 \renewcommand{\algorithmicensure}{\textbf{Output:}}
 \REQUIRE Stage 1 data $(x_i, z_i)$, Stage 2 data $(\tilde y_i \tilde z_i)$, Regularization parameters $(\lambda_1, \lambda_2)$. Initial values $\hat{\theta}_X,\hat{\theta}_Z$. Mini-batch size $(m_b, n_b)$. Number of updates in each stage $(T_1, T_2)$.
 \ENSURE Estimated structural function $\hat{f}_\mathrm{struct}(x)$
 \REPEAT
   \STATE Sample $m_b$ stage 1 data $(x^{(b)}_i, z^{(b)}_i)$ and  $n_b$ stage 2 data $(\tilde y^{(b)}_i, \tilde z^{(b)}_i)$.
  \FOR{$t=1$ to $T_1$}
    \STATE Return function $\hat{\vec{V}}^{(m_b)}(\hat{\theta}_X,\theta_Z)$ in \eqref{eq:stage1-sol}  using $(x^{(b)}_i, z^{(b)}_i)$
    \STATE Update $\hat{\theta}_Z \leftarrow \hat{\theta}_Z - \alpha \nabla_{\theta_Z} \mathcal{L}^{(m_b)}_1(\hat{\vec{V}}^{(m_b)}(\hat{\theta}_X,\theta_Z), \theta_Z)|_{\theta_Z=\hat{\theta}_Z}$ \hspace{5mm}\textit{$\backslash\backslash$ Stage 1 learning}
  \ENDFOR
  \FOR{$t=1$ to $T_2$}
    \STATE Return function $\hat{\vec{u}}^{(n_b)}(\theta_X,\hat{\theta}_Z)$ in  \eqref{eq:stage2-sol} using $ (\tilde y^{(b)}_i, \tilde z^{(b)}_i)$ and function $\hat{\vec{V}}^{(m_b)}(\theta_X,\hat\theta_Z)$\!
    \STATE Update $\hat{\theta}_X \leftarrow \hat{\theta}_X -\alpha\nabla_{\theta_X} \mathcal{L}^{(n_b)}_2(\hat{\vec{u}}^{(n_b)}(\theta_X,\hat{\theta}_Z), \theta_X)|_{\theta_X=\hat{\theta}_X}$ \hspace{5mm}\textit{$\backslash\backslash$ Stage 2 learning}
  \ENDFOR
  \UNTIL{\textbf{convergence}}
  \STATE Compute $\hat{\vec{u}}^{(n)}:=\hat{\vec{u}}^{(n)}(\hat{\theta}_X,\hat{\theta}_Z)$ from \eqref{eq:stage2-sol} using entire dataset.
 \RETURN $\hat{f}_\mathrm{struct}(x) = (\hat{\vec{u}}^{(n)})^\top \vec{\psi}_{\hat{\theta}_X}(x)$
 \end{algorithmic} 
 \end{algorithm}


In Appendix~\ref{sec:consistency}, we provide regularity conditions under which the function learned by DFIV converges to the true structural function in probability. The derivation is based on Rademacher complexity bounds \citep{FoundationOfML}.%, which is applicable not only to neural networks but also to various machine learning techniques.


